\chapter{Mathematical Reference}
\label{app:math-reference}

\section{Algebra}

For $ax+b=0$, the solution is $x=-b/a$, provided $a\neq0$. A linear function $y=a+bx$ has vertical intercept $a$ and slope $b$. The slope between two points is $\Delta y/\Delta x$.

\section{Derivatives}

\[
\frac{d}{dx}c=0,\qquad
\frac{d}{dx}x^n=nx^{n-1},\qquad
\frac{d}{dx}[f(x)+g(x)]=f'(x)+g'(x).
\]

For functions of several variables, a partial derivative changes one variable while holding the others constant:
\[
MU_x=\frac{\partial u}{\partial x},
\qquad
MP_L=\frac{\partial q}{\partial L}.
\]

\section{Optimisation checklist}

\begin{enumerate}
  \item State the objective and constraint.
  \item Substitute to reduce the number of choice variables when appropriate.
  \item Calculate the first-order condition.
  \item Check feasibility and, where relevant, the second-order condition or endpoints.
  \item Interpret the answer economically, including units and assumptions.
\end{enumerate}

\chapter{Selected Answers to Exercises}
\label{app:answers}

These answers provide a check on method and result. They are not substitutes for showing the reasoning required in the exercises.

\section{Chapters 1--5}

\begin{enumerate}
  \item[Chapter 1, Graphical Exercise 6.] The feasible set is $3x+2y\leq120$, and the frontier is $y=60-\frac{3}{2}x$. The intercepts are $(40,0)$ and $(0,60)$. Point $A$ is efficient; $B$ is feasible but inefficient; and $C$ is infeasible. The opportunity cost of one additional economics session is $1.5$ political-philosophy sessions. For Analytical Exercise 8, Marta's opportunity cost is \euro36.
  \item[Chapter 2, Analytical Exercise 8.] The solutions are $x=60$, $x=20$, and $x=4$ or $x=-6$. For Optimisation Exercise 12, $B'(h)=18-2h$, so $h^*=9$; $B''(h)=-2<0$, and $B(9)=81$. For Optimisation Exercise 13,
  \[
  \pi(q)=R(q)-C(q)=48q-2q^2-10.
  \]
  Hence $\pi'(q)=48-4q=0$ gives $q^*=12$, $\pi''(q)=-4<0$, and maximum profit is $\pi(12)=278$.
  \item[Chapter 3, Graphical Exercise 6.] Strong monotonicity gives $B\succ A$, $C\succ A$, and $D\succ A$. Neither $B$ nor $C$ has at least as much of both goods as the other, so strong monotonicity alone cannot rank them. For Analytical Exercise 9, transitivity gives $A\succeq D$. If $D\succeq A$, then $B\succeq C\succeq D\succeq A$, contradicting $A\succ B$. Hence $A\succ D$. For Analytical Exercise 13, use $A\sim B$ and $B\sim C$ to obtain $A\succeq B\succeq C$ and $C\succeq B\succeq A$; thus $A\sim C$.
  \item[Chapter 4, Graphical Exercise 6.] The bundles $(2,3)$, $(3,4)$, and $(6,3)$ lie respectively on $u=6$, $u=12$, and $u=18$. For Analytical Exercise 9,
  \[
  MU_x=\frac{1}{3}x^{-2/3}y^{2/3},
  \qquad
  MU_y=\frac{2}{3}x^{1/3}y^{-1/3},
  \qquad
  MRS_{x,y}=\frac{y}{2x}.
  \]
  For Analytical Exercise 10, the utility level is $18$, so $x(y+2)=18$, or $y=18/x-2$; with non-negative consumption, $0<x\leq9$. For Analytical Exercise 11, the L-shaped curve has kink $(5,10)$ and consists of $x=5$, $y\geq10$, together with $y=10$, $x\geq5$. For Analytical Exercise 12, $e^z$ is strictly increasing, so $e^{xy}$ represents the same preferences as $xy$; both MRS values are $y/x$. For Analytical Exercise 13, the MRS values are $6$, $3/2$, and $2/3$, respectively.
  \item[Chapter 5, Exercise 9.] The optimum is $(x,y)=(5,10)$.
\end{enumerate}

\section{Chapters 6--10}

\begin{enumerate}
  \item[Chapter 6, Exercise 7.] Where both consumers demand positive amounts, market demand is $Q=26-3p$.
  \item[Chapter 7, Exercise 7.] At $p=10$, elasticity is $-0.5$; at $p=25$, it is $-2$; at $p=40$, it is $-4$.
  \item[Chapter 8, Exercise 8.] $MP_L=K^{1/2}L^{-1/2}$, $MP_K=L^{1/2}K^{-1/2}$, and $MRTS_{L,K}=K/L$.
  \item[Chapter 9, Exercise 8.] $MR=80-4q$, $MC=8$, so the profit-maximising output is $q=18$.
  \item[Chapter 10, Exercise 7.] $FC=100$, $VC=4q+q^2$, $AVC=4+q$, $ATC=100/q+4+q$, and $MC=4+2q$.
\end{enumerate}

\section{Chapters 11--15}

\begin{enumerate}
  \item[Chapter 11, Exercise 7.] Equilibrium is $P=30$ and $Q=40$.
  \item[Chapter 12, Exercise 7.] $MR=100-4Q=20$ gives $Q_M=20$, $P_M=60$, and profit $=60(20)-[10+20(20)]=790$.
  \item[Chapter 13, Exercise 7.] The symmetric Cournot equilibrium is $q_1=q_2=38/3$.
  \item[Chapter 14, Exercise 7.] Efficiency requires $MRS^A_{x,y}=MRS^B_{x,y}$, together with feasibility.
  \item[Chapter 15, Exercise 5.] Pollution taxes use prices to decentralise abatement; direct regulation can be simpler where monitoring emissions is difficult; no policy leaves the external cost unaddressed.
\end{enumerate}

\section{Chapters 16--20}

\begin{enumerate}
  \item[Chapter 16, Analytical Exercise 8.] Market output is $80$; efficient output is $60$; the corrective tax is \euro20 per unit.
  \item[Chapter 17, Analytical Exercise 9.] Both players have dominant strategy $A2/B2$ respectively; $(A2,B2)$ is the Nash equilibrium and is Pareto dominated by $(A1,B1)$.
  \item[Chapter 18, Analytical Exercise 7.] The student who values the place at \euro30 receives it and pays \euro20.
  \item[Chapter 19, Analytical Exercise 5.] Truthful bidding is weakly dominant because a bid affects winning but not the price paid conditional on winning.
  \item[Chapter 20, Analytical Exercise 5.] An opt-out pension default changes participation by making saving the path of least resistance while allowing individuals to opt out.
\end{enumerate}

\backmatter

\chapter{Glossary}

\begin{description}
  \item[Allocation] Distribution of scarce resources among people or uses.
  \item[Adverse selection] Hidden information before a contract that distorts participation.
  \item[Assumption] Simplifying statement used to isolate a relationship in a model.
  \item[Average cost] Total cost per unit of output.
  \item[Average rate of change] Change in an output divided by the corresponding change in an input over an interval.
  \item[Axis] A reference line in a graph used to measure and locate values of a variable.
  \item[Budget constraint] The combinations of goods affordable with given income and prices.
  \item[Candidate optimum] A feasible point identified as a possible best choice that still requires checking.
  \item[Cardinal utility] Utility with numerical magnitudes that have a role within a specified model and restricted transformations.
  \item[Ceteris paribus] With other relevant factors held constant.
  \item[Choice set] The alternatives available for an agent to choose from.
  \item[Cobb--Douglas preferences] Preferences represented by a utility function of the form $x^\alpha y^\beta$ with positive exponents.
  \item[Completeness] Property of preferences under which every pair of alternatives can be compared.
  \item[Constraint] Limit on the feasible choices of an economic agent.
  \item[Consumption bundle] Ordered list of quantities of goods consumed by an individual or household.
  \item[Consumer surplus] Willingness to pay minus expenditure.
  \item[Continuity of preferences] Regularity property ruling out arbitrary jumps in preference rankings.
  \item[Convex preferences] Preferences for which mixtures of equally valued bundles are at least as good as the endpoints.
  \item[Dependent variable] A variable represented as an output of a function.
  \item[Derivative] Instantaneous rate of change of a function.
  \item[Diminishing marginal rate of substitution] Decline in MRS as a consumer moves along an indifference curve toward more of one good.
  \item[Dominance] A relation in which one bundle contains at least as much of every desirable good and more of at least one.
  \item[Domain] The values for which a function is mathematically and economically relevant.
  \item[Economic agent] Individual, household, firm, or institution that makes choices.
  \item[Economic model] Simplified representation used to explain or predict economic outcomes.
  \item[Elasticity] Percentage responsiveness of one variable to another.
  \item[Endogenous variable] Variable determined within the model being studied.
  \item[Exogenous variable] Variable taken as given by the model being studied.
  \item[Externality] Unpriced effect of an action on a third party.
  \item[Feasible set] All alternatives available under the relevant constraints.
  \item[First-order condition] A marginal condition, often set equal to zero, used to identify a candidate optimum.
  \item[Function] A rule that maps one or more inputs to an output.
  \item[Horizontal intercept] The point where a graph crosses the horizontal axis.
  \item[Independent variable] A variable treated as an input in a function or graph.
  \item[Incentive] Change in costs or benefits that affects behaviour.
  \item[Indifference] Relation in which an agent regards two alternatives as equally good.
  \item[Indifference curve] Set of bundles that provide the same utility to a consumer.
  \item[Indifference map] A family of indifference curves representing a consumer's preferences.
  \item[Interpersonal comparison of utility] Comparison of well-being between people, which requires ethical assumptions beyond ordinal utility.
  \item[Institution] Formal or informal rule structuring social interaction.
  \item[Local maximum] A point whose function value exceeds those at all sufficiently nearby feasible points.
  \item[Marginal analysis] Comparison of the additional benefit and cost of a small change.
  \item[Marginal benefit] Additional benefit from one more unit of an activity.
  \item[Marginal cost] Additional cost of one more unit of output.
  \item[Marginal rate of substitution] Rate at which a consumer trades one good for another while preserving utility.
  \item[Marginal utility] Change in utility associated with a small increase in one good, holding other goods fixed.
  \item[Market equilibrium] Price and quantity at which quantity demanded equals quantity supplied.
  \item[Mechanism design] Design of rules that induce strategic participants to produce desirable outcomes.
  \item[Monopoly] Single seller protected by barriers to entry.
  \item[Monotonicity] Assumption that, holding other goods fixed, more of a desirable good is preferred to less.
  \item[Nash equilibrium] Mutual best responses in a game.
  \item[Non-satiation] Assumption that more of at least one desirable good can improve a consumer's bundle.
  \item[Normative claim] Claim about what ought to be done, involving values or ethical judgement.
  \item[Objective function] Mathematical expression describing the outcome an agent seeks to maximise or minimise.
  \item[Ordinal utility] Utility representation in which only the ranking of bundles, not numerical differences, matters.
  \item[Opportunity cost] Value of the best alternative forgone.
  \item[Parameter] Quantity treated as fixed within a particular model or problem.
  \item[Pareto efficiency] State in which no Pareto improvement is possible.
  \item[Perfect complements] Goods valued in fixed proportions, represented by L-shaped indifference curves.
  \item[Perfect substitutes] Goods that can replace one another at a constant rate, represented by straight indifference curves.
  \item[Partial derivative] Rate of change with respect to one variable while the other arguments are held fixed.
  \item[Price discrimination] Charging different buyers different prices not explained by cost.
  \item[Positive claim] Claim about how the world works that can be assessed with theory and evidence.
  \item[Preference relation] A ranking of alternatives according to an agent's judgements.
  \item[Property rights] Rules identifying who may use, control, or transfer a resource.
  \item[Public good] Non-rival and non-excludable good.
  \item[Quasilinear preferences] Preferences represented by $u(x,y)=v(x)+y$, with constant marginal utility of $y$.
  \item[Rational-choice assumption] Modelling assumption that agents pursue objectives within their constraints.
  \item[Rational preferences] Complete and transitive preferences, often with further regularity assumptions for particular applications.
  \item[Reflexivity] Property that every alternative is at least as good as itself.
  \item[Scarcity] Limited resources relative to wants or possible uses.
  \item[Second-order condition] A condition using a second derivative to classify a candidate optimum.
  \item[Slope] Change in the vertical-axis variable associated with a one-unit change in the horizontal-axis variable.
  \item[Critical point] A point at which the first derivative is zero or does not exist.
  \item[Stationary point] A point at which the first derivative is zero.
  \item[Strong monotonicity] Assumption that a bundle with at least as much of every desirable good and more of at least one is strictly preferred.
  \item[Strict convexity] Property under which a genuine mixture of distinct indifferent bundles is strictly preferred to either endpoint.
  \item[Strictly increasing transformation] Transformation that preserves the ordering of every pair of utility values.
  \item[Strict preference] Relation in which one alternative is at least as good as another but the reverse weak preference does not hold.
  \item[Trade-off] The sacrifice of one valuable option to obtain more of another.
  \item[Transitivity] Property that if $A\succeq B$ and $B\succeq C$, then $A\succeq C$.
  \item[Utility] Numerical representation of preferences.
  \item[Utility function] Function assigning numbers to bundles in a way that represents preferences.
  \item[Utility representation] A utility function that preserves an agent's preference ranking.
  \item[Variable] Quantity that can take different values.
  \item[Vertical intercept] The point where a graph crosses the vertical axis.
  \item[Weak preference] Relation $A\succeq B$, meaning that $A$ is at least as good as $B$.
\end{description}

\chapter{References and Further Reading}

\begin{thebibliography}{9}
  \bibitem{varian} Varian, H. R. \emph{Intermediate Microeconomics: A Modern Approach}. W. W. Norton.
  \bibitem{perloff} Perloff, J. M. \emph{Microeconomics}. Pearson.
  \bibitem{pindyck} Pindyck, R. S., and Rubinfeld, D. L. \emph{Microeconomics}. Pearson.
  \bibitem{mascolell} Mas-Colell, A., Whinston, M. D., and Green, J. R. \emph{Microeconomic Theory}. Oxford University Press.
  \bibitem{kahneman} Kahneman, D. \emph{Thinking, Fast and Slow}. Farrar, Straus and Giroux.
  \bibitem{coase} Coase, R. H. ``The Problem of Social Cost.'' \emph{Journal of Law and Economics}, 3, 1960.
\end{thebibliography}
